\section{From Agent Traces to Trust: A Survey of Evidence Tracing and Execution Provenance in LLM Agents}

Yiqi Wang, Jiaqi Zhang, Zhangkai Wu, Taotao Cai, Zirui Liu, Zequn Sun, Mingkai Zheng,
Manqing Dong, Yu-Xuan Qiu, Xuefei Yin, Yanming Zhu. \emph{From Agent Traces to Trust:
A Survey of Evidence Tracing and Execution Provenance in LLM Agents.} arXiv:2606.04990v6.

\cite{wang2026survey}

\subsection{Ключевые идеи, математическое описание метода}
Обзор вводит два понятия и формализует связь между ними. \emph{Execution provenance}~--
типизированный граф прогона; \emph{evidence tracing}~-- его проекция на отношения
поддержки. Лог отвечает на вопрос <<что произошло и в каком порядке>>, граф~-- <<что от чего
зависело>>; порядок во времени связью не считается.

Граф одного прогона (прил.~C.1):
\begin{equation}
G = (V, E, \tau, \lambda, \alpha), \qquad V = E_v \uplus X \uplus A,
\end{equation}
где $\tau$ приписывает узлу тип и делит узлы на единицы-свидетельства $E_v$ (артефакты
с содержимым $\rho(v)$: запрос, документ, вывод инструмента, параметр, утверждение),
единицы выполнения $X$ (события с исполнителем и интервалом времени: шаг рассуждения,
вызов инструмента) и исполнителей $A$; $E$~-- рёбра мультиграфа с отображениями
$\mathrm{src},\mathrm{tgt}$; $\lambda$~-- метка из
$R^+=R\cup\{\text{Support-candidate}\}$, где $R$ содержит девять основных отношений;
$\alpha$~-- атрибуты,
в том числе \emph{происхождение} ребра. Отношения делятся на семейства
\begin{equation}
\begin{aligned}
R_{sem}&=\{\text{Support},\text{Contradict},\text{Invalidate}\},\\
R_{der}&=\{\text{Derive}\},\\
R_{proc}&=\{\text{Use},\text{Generate},\text{Depend-on},\text{Trigger},\text{Update}\}.
\end{aligned}
\end{equation}
с сигнатурами по типам узлов и временными ограничениями (табл.~11): например,
$\text{Derive}: E_v\to E_v$, $t(s)<t(t)$; $\text{Depend-on}: X\to E_v\cup X$, устанавливается
контрфактически. Происхождение ребра: \emph{Instrumented} (записано харнессом),
\emph{Inferred} (вынесено моделью или разметчиком), \emph{Declared} (заявлено агентом),
\emph{Derived} (вычислено из графа). Требование обзора: Inferred и Declared рёбра должны
получать проверку, соответствующую смыслу отношения: например, разметкой поддержки
или вмешательствами для гипотезы о влиянии. Сам факт записи не удостоверяет связь.

\subsection{Алгоритм}
Проекция $\pi: G\mapsto\pi(G)$ (прил.~C.1), пять правил, вычисляемых на исходном $G$:
\begin{enumerate}
  \item сохранить рёбра с $\lambda(e)\in R_{sem}\cup R_{der}\cup\{\text{Support-candidate}\}$;
  \item отбросить все рёбра из $R_{proc}$;
  \item добавить $\text{Support-candidate}(d,c)$ для $d\in E_v$, $c\in C$ (утверждающие
        единицы), если в $G$ есть путь потока содержимого из $d$ в $c$ (Derive вперёд,
        Generate вперёд, Use назад, хотя бы одно ребро Use или Generate) и между $d$ и $c$
        ещё нет Support или Invalidate; такое ребро утверждает возможный вклад содержимого,
        но не поддержку, его происхождение~-- Derived;
  \item оставить узлы-свидетельства, инцидентные сохранённым рёбрам, и шаги~-- цели Support
        или Invalidate;
  \item сохранить атрибуты содержимого, хешей, времени, доверия, происхождения и assurance;
        для Support-candidate сохранить идентификатор пути-свидетеля;
        отбросить атрибуты выполнения.
\end{enumerate}
Свойства: $\pi$ не инъективна (разные прогоны дают одну трассировку свидетельств), обоснована
(не добавляет Support), отбрасывает всё влияние без поддержки (Depend-on, Trigger),
идемпотентна: $\pi(\pi(G))=\pi(G)$.

\subsection{Наборы данных}
Собственных данных нет. Каталог бенчмарков (табл.~8, разд.~6): RAG-атрибуция (ALCE, FActScore,
RAGTruth), трассы и безопасность инструментов (TRAIL, MAST, AgenTracer/TracerTraj, InjecAgent,
AgentDojo, ToolEmu, OpenAgentSafety, MCP-SafetyBench), память и многоагентность.
Вывод разд.~6.5 ограничен собственным каталогом: ни один из его 25 пунктов не оценивает
полную цепочку провенанса. Отдельные отношения уже размечаются: Support в атрибуции,
source-to-sink в taint-бенчмарках; единого эталона полного графа прогона нет.

\subsection{Метрики}
Табл.~7: четыре группы со статусом зрелости. Established~-- evidence recall, citation precision,
faithfulness (из RAG-атрибуции). Partly established~-- локализация отказа (TRAIL/MAST).
\emph{Proposed}~-- trace completeness, provenance accuracy, dependency coverage,
temporal consistency: общепринятых протоколов пока нет. Разд.~6.4 предлагает эталон
инструментированного харнесса для completeness, ablation/replay для dependency coverage
и подсчёт временных нарушений для temporal consistency. Для accuracy обсуждаются
допустимое множество графов, взвешенные оценки рёбер или канонический граф.
Это программа оценки, а не проведённый авторами эксперимент восстановления связей.

\subsection{Модели}
Не применимо (обзор). Систематизированы защиты на информационных потоках (CaMeL, FIDES,
NeuroTaint, Agent-Sentry; разд.~4.2--4.3) и системы наблюдаемости (AgentOps, AgentTrace,
PROV-AGENT; табл.~3).
